\documentclass[a4paper,12pt]{article}
\usepackage[finnish]{babel}
\usepackage[utf8]{inputenc}
\usepackage[T1]{fontenc}
\usepackage{lmodern}
\usepackage{eurosym}
\usepackage[pdftex,colorlinks]{hyperref}
\renewcommand{\thesubsection}{\arabic{subsection}}
\title{Etelä-Espoon sosialidemokraattinen työväenyhdistys r.y.}
\author{EED}
\date{21.01.2026}
\begin{document}
\maketitle
\tableofcontents
\section*{Johtokunnan kokous}
\subsection*{Paikka ja aika}
Tapiola Garden, keskiviikkona 21.01.2026 kello 18:00.
\section*{Esityslista}
\subsection{Kokouksen avaus}
\subsection{Kokouksen laillisuus ja päätösvaltaisuus}
\subsection{Kokouksen ääntenlaskijoiden valinta}
\subsection{Johtokunnan järjestäytyminen}
Valitaan johtokunnalle sihteeri, rahastonhoitaja, jäsenasiainhoitaja ja tiedotusvastaava.
\subsection{Taloustilanne}
Päätetään Falklammen mökin hinnaston muutoksista. Hinnat riippuvat siitä, onko vuokraaja jäsen, vuokrataanko mökki koko viikoksi vai viikonlopuksi ja onko vuokra-aika kesäsesonkina. Hinnoissa on merkittävää korotuspainetta, johtuen mökin vuosittaisista kunnostustarpeista ja siitä seikasta, että hinnat ovat jääneet vuosien saatossa huomattavasti jälkeen alueen yleisestä vuokramökkien hintatasosta. Puheenjohtaja Markku Sistonen esittelee asian tarkemmin.
\subsection{Jäsenasiat}
\subsection{Puoluekokousedustajat}
Keskustellaan puoluekokousedustajien valinnasta. Sähköposteissa Helena Tiitinen ja Maria Guzenina ovat ilmaisseet kiinnostustaan puoluskokousedustajaksi. Yhdistys joutuu järjestämään tämän vuoksi ylimääräisen jäsenkokouksen.
\subsection{Eduskuntavaaliehdokkaat}
Keskustellaan eduskuntavaaliehdokkaiden nimeämistestä. Sähköposteissa Helena Tiitinen on ilmaissut kiinnostustaan ehdokkuuteen.
\subsection{Kevään tapahtumat}
\begin{itemize}
\item{Kreikan matka}
\end{itemize}
\subsection{Muut asiat}
\subsection{Kokouksen päättäminen}
\end{document}
